\section{Definitions}
\begin{definition}[Code]
\label{def:Code}
\lean{Code}
\leanok
A code $C$ of block length $n$ over an alphabet $\alpha$ is a subset of $\alpha^{n}$.
Typically, we will use $q$ to denote the alphabet size $|\Sigma|$.
\end{definition}

\begin{definition}[Alphabet Size]
    \label{def:q}
    \lean{Code.q}
    \leanok
    The alphabet size is $q := \lvert\alpha\rvert$, the cardinality of the finite alphabet $\alpha$.
\end{definition}

\begin{definition}[Dimension]
    \label{def:dim}
    \lean{Code.dim}
    \leanok
    \uses{def:Code, def:q}
    The dimension $k$ of a code $C$ is
    \[
        k := \log_{q}(\lvert C\rvert),
    \]
    where $q$ is the alphabet size and $\lvert C\rvert$ is the number of codewords.
\end{definition}
 
\begin{definition}[Rate]
    \label{def:rate}
    \lean{Code.rate}
    \leanok
    \uses{def:Code, def:dim}
    The rate $R$ of a code $C$ is
    \[
        R := \frac{k}{n}.
    \]
    where $k$ is the dimension and block length of $n$.
\end{definition}
